\scanpage{352}
\mspage{91}
% archive sheet number: 351
alors ramené au cas où $\boxed{\unsure{I|\tau}=\Delta_\varepsilon^2}$\inserted{, $\tau'=\Delta_\varepsilon^2$}. Sinon, au cas \unsure{où|en} $\boxed{\tau=\Delta_\varepsilon\Delta_{-\varepsilon}^{\ell_2}\Delta_\varepsilon}$, \inserted{$\tau'=\Delta_\varepsilon^2$} et où \deleted{\unsure{une}} les deux flèches critiques sont encore la première et la dernière, et $\tau'=\Delta_\varepsilon^2$. Ce sont là les deux cas irréductibles, qu'il \deleted{\unsure{y}} faudra traiter \deleted{\unsure{de façon}} à part.

\deleted{2)} Ces cas admis, on a donc moyen pour déformer $F$ de façon à satisfaire a). \deleted{Soient} Ensuite, on veut

2) Assurer que la suite des intervalles critiques $I_{j_1},I_{j_2}\ldots I_{j_{r'}}$ qui \unsure{contient|contiennent} les flèches critiques $\underline{u}_{i_1},\ldots,\underline{u}_{i_{r'}}$, \deleted{\unsure{soit}} soit \unsure{de|la} suite $I_1,\ldots I_{r'}$.
C'est encore \unsure{immédiat$^{\mathrm{t}}$|immédiat} par l'hyp\unsure{.} de récurrence, le seul cas irréd. étant celui où $r'=1$\inserted{ (une seule flèche critique)}, et où il \deleted{\unsure{n'a que deux intervalles critiques, et}} s'agit de faire descendre cette flèche critique jusqu'à $\underline{u}_1$. Si $I_1$ est de longueur $>1$, alors on fait descendre en $\underline{u}_2$ par hyp. de récurrence, et le cas irréductible est celui où $\boxed{\tau=\Delta_\varepsilon^2}$\inserted{, $\tau'=\Delta_\varepsilon$} une flèche critique $\underline{u}_2$, la faire descendre en $\underline{u}_1$.

Si par contre $I_1$ de longueur $1$, le cas \unsure{irréductible} est $\boxed{\tau=\Delta_\varepsilon\Delta_{-\varepsilon}^{\ell_2}\Delta_\varepsilon}$, \deleted{la flèche} $\tau'=\Delta_\varepsilon$, une flèche critique $\underline{u}_{\ell_2+2}$, la faire descendre en $\underline{u}_1$.

Sous réserve de traiter ces \inserted{deux} nouveaux cas irréductibles, on peut déformer $F$ en un foncteur qui satisfait à a) et à c). Reste à normaliser chacun des $F_\alpha:I_\alpha\mathrel{\unsure{-|\to}}I'_\alpha$ ($\alpha\in[1,r']$).

Donc on est ramené à $\Delta_\varepsilon^{\unsure{n|N}}\longrightarrow\Delta_\varepsilon$~: \unsure{$\{$}une seule flèche critique\unsure{$\}$} $\underline{u}_\alpha$, $\alpha\in[1,n]$ $\}$, à faire descendre en $\underline{u}_{\deleted{\unsure{2}}\,\unsure{1}}$. Si $n\geq2$, par hyp. de récurrence elle descend \unsure{en} $\underline{u}_2$, cas supposé traité en $2^{\mathrm{o}}$). Donc
